\documentclass[11pt,a4paper]{article}
\usepackage[a4paper,margin=2.25cm]{geometry}
\usepackage[spanish,es-noquoting]{babel}
\usepackage[T1]{fontenc}
\usepackage[utf8]{inputenc}
\usepackage{lmodern,microtype}
\usepackage{amsmath,amssymb,amsthm,mathtools}
\usepackage{booktabs,array,longtable,enumitem}
\usepackage{xcolor}
\usepackage[hidelinks]{hyperref}

\pagecolor{white}
\color{black}
\setlength{\parindent}{0pt}
\setlength{\parskip}{0.65em}

\newtheorem{teorema}{Teorema}
\newtheorem{proposicion}[teorema]{Proposici\'on}
\newtheorem{definicion}[teorema]{Definici\'on}
\newtheorem{observacion}[teorema]{Observaci\'on}
\newtheorem{corolario}[teorema]{Corolario}

\newcommand{\dr}{\operatorname{dr}_9}
\newcommand{\APP}{\mathrm{APP}}
\newcommand{\Bc}{\mathcal B_{\rm c}}
\newcommand{\Bang}{\mathcal B_{\rm ang}}
\newcommand{\Tang}{\mathcal T_{\rm ang}}
\newcommand{\HMT}{\mathrm{HMT}}
\newcommand{\MD}{\mathrm{MD}}
\newcommand{\R}{\mathbb R}
\newcommand{\Z}{\mathbb Z}
\newcommand{\F}{\mathbb F}
\newcommand{\Id}{I}
\newcommand{\rad}{\mathfrak m}
\newcommand{\DeltaAPP}{\Delta_{\rm APP}}
\newcommand{\hbarbar}{\bar\hbar}
\newcommand{\Spec}{\operatorname{Spec}}
\newcommand{\tr}{\operatorname{tr}}
\newcommand{\sgn}{\operatorname{sgn}}

\title{\bfseries Ocho teoremas del n\'ucleo APP\\
\large Operador central two-way, geometr\'ia Lorentz, interfaz angular y caracteres de masa}
\author{Ventana t\'ecnica HMT}
\date{}

\begin{document}
\maketitle

\begin{abstract}
Este texto formaliza el n\'ucleo multiplicativo de la APP como un objeto algebraico central de HMT. El bloque
\[
\Bc=\begin{pmatrix}7&2\\2&7\end{pmatrix}
\]
se interpreta como operador two-way, elemento del \`algebra split \(\R[R]/(R^2-I)\), compresi\'on local del comportamiento dimensional de las hojas aditiva y multiplicativa, generador de una geometr\'ia de Lorentz en dimensi\'on \(1+1\), y puente hacia la Mec\'anica Dimensional mediante el levantamiento angular
\[
\Bang=A I+C^\ast R.
\]
El texto organiza ocho resultados: la familia tr\'itica de bloques two-way; la compresi\'on Hensel de la APP; la localizaci\'on radical del defecto \(54\); la normalizaci\'on Lorentz del centro; la compresi\'on global de la red \(3,6,9\); la interfaz angular \((A,C^\ast)\); la ley de masas como car\'acter de los canales espectrales; y la extracci\'on nuclear de las estructuras \(6/8\), \(-1/12\) y Moonshine. Se separan pruebas algebraicas, lecturas HMT y puntos que requieren m\'odulos posteriores.
\end{abstract}

\tableofcontents
\newpage

\section{Contrato editorial: definici\'on afirmativa}

El n\'ucleo APP se expone como sistema algebraico tipado. Cada entidad se define por su soporte, su ley de composici\'on, sus invariantes y su funci\'on dentro del corpus. La comparaci\'on con formulaciones exteriores queda reservada a los apartados de lectura o auditor\'ia. En el cuerpo principal, cada objeto se presenta por lo que es y por lo que hace.

El prop\'osito de esta ventana es fijar el puente:
\[
\text{APP}\longrightarrow \text{\'algebra split/Lorentz}\longrightarrow \text{interfaz angular}\longrightarrow \text{Mec\'anica Dimensional}.
\]

El resultado principal puede enunciarse de forma compacta:
\[
\boxed{\text{el centro de APP es la bisagra espectral entre aritm\'etica, Lorentz y materia.}}
\]

\section{Definiciones b\'asicas}

\begin{definicion}[ra\'iz digital de orden nueve]
Para un entero positivo \(n\), escribimos
\[
\dr(n)=1+((n-1)\bmod 9).
\]
As\'i, el residuo cero m\'odulo nueve se representa por el d\'igito \(9\).
\end{definicion}

\begin{definicion}[hojas APP]
La hoja multiplicativa de la APP de orden nueve es la matriz
\[
M_{ij}=\dr(ij),\qquad 1\le i,j\le 9.
\]
La hoja aditiva de transporte se toma como
\[
A_{ij}=\dr(i+j-1),\qquad 1\le i,j\le 9.
\]
La convenci\'on alternativa \(\dr(i+j)\) s\'olo permuta c\'iclicamente los bloques de la compresi\'on aditiva y conserva las sumas globales.
\end{definicion}

\begin{definicion}[involuci\'on two-way]
Definimos
\[
R=\begin{pmatrix}0&1\\1&0\end{pmatrix},\qquad R^2=I.
\]
Sus proyectores son
\[
P_+=\frac{I+R}{2},\qquad P_-=\frac{I-R}{2}.
\]
\end{definicion}

\section{Teorema I: familia tr\'itica de bloques two-way}

\begin{teorema}[familia tr\'itica de pares opuestos]
En la hoja multiplicativa de APP, los pares opuestos de unidades
\[
\{1,8\},\qquad \{2,7\},\qquad \{4,5\}
\]
producen tres bloques two-way de la forma
\[
\mathcal B_\tau=(4+3\tau)I+(5-3\tau)R,
\qquad \tau\in\{-1,0,+1\}.
\]
Su espectro es
\[
\boxed{\Spec(\mathcal B_\tau)=\{9,\,6\tau-1\}.}
\]
En consecuencia,
\[
\boxed{\tau=\frac{\lambda_-(\tau)+1}{6}.}
\]
El trit se recupera como coordenada normalizada del canal diferencial.
\end{teorema}

\begin{proof}
Los tres pares opuestos dan, respectivamente,
\[
\mathcal B_{-1}=\begin{pmatrix}1&8\\8&1\end{pmatrix},\quad
\mathcal B_{0}=\begin{pmatrix}4&5\\5&4\end{pmatrix},\quad
\mathcal B_{+1}=\begin{pmatrix}7&2\\2&7\end{pmatrix}.
\]
Todos tienen suma de fila \(9\), de modo que el vector \((1,1)\) es propio con autovalor \(9\). El vector \((1,-1)\) tiene autovalor \(a-b\) para un bloque \(aI+bR\). En los tres casos,
\[
(a,b)=(1,8),(4,5),(7,2),
\]
por lo que
\[
a-b=-7,-1,5.
\]
Estos tres valores son \(6\tau-1\) para \(\tau=-1,0,+1\). La f\'ormula matricial
\[
\mathcal B_\tau=(4+3\tau)I+(5-3\tau)R
\]
reproduce exactamente los tres casos. La f\'ormula del trit se obtiene despejando \(\tau\).
\end{proof}

\begin{observacion}[lectura HMT]
El canal com\'un se cierra siempre en \(9\). El canal diferencial contiene un vac\'io de referencia \(-1\) y una activaci\'on hex\'adica \(6\tau\). La familia de bloques realiza espectralmente la unidad tr\'itica.
\end{observacion}

\section{Teorema II: operador central y compresi\'on espectral}

\begin{teorema}[operador central two-way]
El bloque central multiplicativo
\[
\Bc=\begin{pmatrix}7&2\\2&7\end{pmatrix}=7I+2R
\]
posee la descomposici\'on espectral
\[
\boxed{\Bc=9P_+ +5P_-.}
\]
Su determinante es
\[
\det\Bc=45,
\]
y su salto espectral es
\[
9-5=4.
\]
\end{teorema}

\begin{proof}
Para una matriz \(aI+bR\), los autovalores en los subespacios generados por \((1,1)\) y \((1,-1)\) son \(a+b\) y \(a-b\). Con \(a=7,b=2\), obtenemos \(9\) y \(5\). La descomposici\'on por proyectores sigue de la identidad
\[
aI+bR=(a+b)P_++(a-b)P_-.
\]
El determinante es \((7+2)(7-2)=9\cdot5=45\). El salto espectral es \(9-5=4\).
\end{proof}

\begin{observacion}[compresi\'on dimensional]
El bloque central realiza localmente dos l\'imites de la APP dimensional: el centro aditivo estable \(5\) y el cierre multiplicativo asint\'otico \(9\). El salto \(4\) codifica el defecto asint\'otico entre esas dos lecturas. Esta afirmaci\'on requiere el m\'odulo dimensional donde se definen las cantidades \(\mu_{+,d}\) y \(\mu_{\times,d}\). En ese m\'odulo, el bloque central act\'ua como compresi\'on local de los l\'imites globales.
\end{observacion}

\section{Teorema III: semigrupo exacto y memoria diferencial}

\begin{teorema}[semigrupo del centro]
Para todo \(n\ge0\),
\[
\boxed{\Bc^n=9^nP_+ +5^nP_-}
\]
y, en la base original,
\[
\boxed{
\Bc^n=\frac12
\begin{pmatrix}
9^n+5^n&9^n-5^n\\
9^n-5^n&9^n+5^n
\end{pmatrix}.}
\]
Adem\'as,
\[
\boxed{9^{-n}\Bc^n=P_+ +\left(\frac59\right)^nP_-\longrightarrow P_+.}
\]
\end{teorema}

\begin{proof}
La primera igualdad se obtiene por c\'alculo funcional sobre los proyectores ortogonales \(P_+\) y \(P_-\). La forma en la base original se obtiene sustituyendo
\[
P_+=\frac12\begin{pmatrix}1&1\\1&1\end{pmatrix},\qquad
P_-=\frac12\begin{pmatrix}1&-1\\-1&1\end{pmatrix}.
\]
La convergencia normalizada es inmediata, pues \((5/9)^n\to0\).
\end{proof}

\begin{observacion}[renormalizaci\'on exacta]
La iteraci\'on privilegia el modo com\'un y conserva el modo diferencial con peso relativo \((5/9)^n\). La cantidad \(5/9\) es la tasa espectral de persistencia de la diferencia two-way.
\end{observacion}

\section{Teorema IV: normalizaci\'on Lorentz-APP}

\begin{teorema}[boost Lorentziano del centro]
La normalizaci\'on
\[
\widehat{\Bc}=\frac{\Bc}{\sqrt{45}}
\]
es un elemento de norma uno del \`algebra split \(\R[R]/(R^2-I)\). M\'as precisamente,
\[
\boxed{\widehat{\Bc}=\exp(\eta_0R),}
\]
con
\[
\boxed{\eta_0=\operatorname{artanh}\frac27=\frac12\log\frac95.}
\]
\end{teorema}

\begin{proof}
Tenemos
\[
\widehat{\Bc}=\frac7{\sqrt{45}}I+\frac2{\sqrt{45}}R.
\]
Como
\[
\left(\frac7{\sqrt{45}}\right)^2-\left(\frac2{\sqrt{45}}\right)^2=1,
\]
existe \(\eta_0\) con
\[
\cosh\eta_0=\frac7{\sqrt{45}},\qquad
\sinh\eta_0=\frac2{\sqrt{45}}.
\]
Dado que \(R^2=I\),
\[
e^{\eta_0R}=\cosh\eta_0 I+\sinh\eta_0 R.
\]
Finalmente,
\[
\tanh\eta_0=\frac{2}{7},
\]
de donde
\[
\eta_0=\operatorname{artanh}\frac27=\frac12\log\frac{1+2/7}{1-2/7}=\frac12\log\frac95.
\]
\end{proof}

\begin{observacion}[significado de Lorentz]
La norma
\[
N(aI+bR)=a^2-b^2
\]
posee firma \((1,1)\). En coordenadas propias \(P_+,P_-\), el operador normalizado multiplica un rayo por \(\sqrt{9/5}\) y el otro por \(\sqrt{5/9}\), conservando el producto. La geometr\'ia causal m\'inima aparece ya en el centro aritm\'etico de APP.
\end{observacion}

\section{Teorema V: compresi\'on Hensel y localizaci\'on del defecto}

\begin{definicion}[clases Hensel]
Agrupamos los d\'igitos por residuo m\'odulo tres:
\[
C_1=\{1,4,7\},\qquad C_2=\{2,5,8\},\qquad C_0=\{3,6,9\}.
\]
La clase \(C_0\) es el radical \(3,6,9\) de \(\Z/9\Z\).
\end{definicion}

\begin{teorema}[compresi\'on por clases]
La suma por bloques de la hoja multiplicativa produce
\[
\boxed{
S_\times=
\begin{pmatrix}
36&45&54\\
45&36&54\\
54&54&81
\end{pmatrix}.}
\]
La matriz de medias es
\[
\boxed{
Q_\times=
\begin{pmatrix}
4&5&6\\
5&4&6\\
6&6&9
\end{pmatrix}.}
\]
Para la hoja aditiva de transporte se obtiene
\[
\boxed{
Q_+=
\begin{pmatrix}
4&5&6\\
5&6&4\\
6&4&5
\end{pmatrix}.}
\]
La diferencia comprimida es
\[
\boxed{
Q_\times-Q_+=
\begin{pmatrix}
0&0&0\\
0&-2&2\\
0&2&4
\end{pmatrix}.}
\]
De ella se extrae la escalera
\[
\boxed{2\longrightarrow 6\longrightarrow54.}
\]
\end{teorema}

\begin{proof}
La matriz \(S_\times\) se obtiene sumando los valores \(\dr(ij)\) con \(i\in C_a,j\in C_b\). Cada bloque tiene nueve entradas, y por eso la matriz de medias es \(Q_\times=S_\times/9\). La hoja aditiva de transporte \(\dr(i+j-1)\) se comprime de la misma forma. La resta es directa. Su traza es \(2\), la suma de sus entradas es \(6\), y al desplegarla en bloques de nueve entradas se obtiene \(9\cdot6=54\).
\end{proof}

\begin{corolario}[soporte radical del defecto]
El defecto global
\[
\DeltaAPP=S_\times(9)-S_+(9)=54
\]
est\'a soportado por el sector no invertible de \(\Z/9\Z\), es decir, por la interacci\'on con el radical \(3,6,9\).
\end{corolario}

\begin{proof}
Las seis filas unitarias de la hoja multiplicativa suman \(45\), igual que las filas aditivas. Las filas \(3\) y \(6\) suman \(54\), y la fila \(9\) suma \(81\). Por tanto,
\[
S_\times=6\cdot45+2\cdot54+81=459,
\]
frente a
\[
S_+=9\cdot45=405.
\]
La diferencia es
\[
2(54-45)+(81-45)=54.
\]
\end{proof}

\section{Teorema VI: red radical y holograf\'ia local-global}

\begin{teorema}[compresi\'on local de la red \(3,6,9\)]
La uni\'on de las filas y columnas \(3,6,9\) de la hoja multiplicativa tiene suma total
\[
297.
\]
Su cuadrado interior suma
\[
81,
\]
y sus brazos exteriores suman
\[
216.
\]
Adem\'as,
\[
216=3\cdot72,
\qquad
81=3\cdot27,
\qquad
297=3\cdot99.
\]
En consecuencia, el par local \((72,27)\) le\'ido en el bloque central es la compresi\'on por factor tres de la descomposici\'on global de la red radical.
\end{teorema}

\begin{proof}
En las intersecciones de las filas y columnas \(3,6,9\), todo producto es m\'ultiplo de nueve y se lee como \(9\), luego las nueve entradas suman \(81\). La suma total de todas las entradas situadas en filas o columnas \(3,6,9\), contando la uni\'on una sola vez, es \(297\). La diferencia \(297-81\) da los brazos exteriores, \(216\). Las identidades con \(72,27,99\) son inmediatas.
\end{proof}

\begin{observacion}[lecturas del centro]
El bloque central permite las lecturas
\[
72+27=99,
\qquad
72-27=45,
\]
\[
77+22=99,
\qquad
77-22=55=54+1.
\]
As\'i, el centro codifica simult\'aneamente el invariante aditivo \(45\), el defecto \(54\), el defecto sellado \(55\) y el cierre \(99\).
\end{observacion}

\section{Teorema VII: interfaz angular y canales f\'isicos}

\begin{definicion}[interfaz angular]
La Mec\'anica Dimensional levanta el operador central a
\[
\boxed{\Bang=AI+C^\ast R,}
\]
donde
\[
\boxed{A=1000\alpha,}
\]
y
\[
\boxed{C^\ast=2(\hbarbar+\alpha).}
\]
Aqu\'i \(\hbarbar\) denota la normalizaci\'on adimensional de la constante reducida de Planck dentro del corpus. La dimensionalizaci\'on f\'isica posterior restituye las unidades metrol\'ogicas.
\end{definicion}

\begin{teorema}[espectro angular]
El operador angular conserva los proyectores del centro APP:
\[
\boxed{\Bang=(A+C^\ast)P_+ +(A-C^\ast)P_-.}
\]
Sus canales propios son
\[
\boxed{\lambda_+^{\rm ang}=A+C^\ast,\qquad \lambda_-^{\rm ang}=A-C^\ast.}
\]
Al sustituir las definiciones de \(A\) y \(C^\ast\), se obtiene
\[
\boxed{A+C^\ast=1002\alpha+2\hbarbar,}
\]
\[
\boxed{A-C^\ast=998\alpha-2\hbarbar.}
\]
\end{teorema}

\begin{proof}
La descomposici\'on espectral de todo operador \(aI+bR\) es \((a+b)P_+ +(a-b)P_-\). Tomando \(a=A,b=C^\ast\) se obtiene el resultado. Las f\'ormulas \(1002\alpha+2\hbarbar\) y \(998\alpha-2\hbarbar\) se obtienen sustituyendo \(A=1000\alpha\) y \(C^\ast=2(\hbarbar+\alpha)\).
\end{proof}

\begin{teorema}[propagador angular]
Sea \(\kappa=\pi/180\). Definimos
\[
\Tang=\exp(-\kappa\Bang).
\]
Entonces
\[
\boxed{\Tang=q_+P_+ +q_-P_-,}
\]
donde
\[
\boxed{q_+=\exp[-\kappa(A+C^\ast)],\qquad q_- =\exp[-\kappa(A-C^\ast)].}
\]
En la base de hojas,
\[
\boxed{\Tang=e^{-\kappa A}\bigl(\cosh(\kappa C^\ast)I-\sinh(\kappa C^\ast)R\bigr).}
\]
Adem\'as,
\[
\boxed{A=-\frac{90}{\pi}\log(q_+q_-),\qquad C^\ast=\frac{90}{\pi}\log\frac{q_-}{q_+}.}
\]
\end{teorema}

\begin{proof}
El c\'alculo funcional sobre los proyectores da los autovalores del exponencial. Como \(R^2=I\),
\[
e^{-\kappa C^\ast R}=\cosh(\kappa C^\ast)I-\sinh(\kappa C^\ast)R.
\]
El producto \(q_+q_-=e^{-2\kappa A}\) y el cociente \(q_-/q_+=e^{2\kappa C^\ast}\) dan las f\'ormulas inversas.
\end{proof}

\begin{observacion}[conjugaci\'on]
La conjugaci\'on part\'icula-antipart\'icula act\'ua como
\[
C^\ast\mapsto-C^\ast,
\qquad
q_+\leftrightarrow q_-.
\]
El modo com\'un \(A\) permanece fijo y la norma split \(A^2-(C^\ast)^2\) permanece invariante.
\end{observacion}

\section{Teorema VIII: masas como caracteres espectrales}

\begin{teorema}[car\'acter de masa]
Sea una especie \(P\) con ledger efectivo \((n_A(P),s_C(P))\). Si
\[
\frac{m(P)}{m_e}=\exp\bigl(n_A(P)A_{\rm rad}+s_C(P)C^\ast_{\rm rad}\bigr),
\]
y si
\[
\lambda_+^{\rm rad}=A_{\rm rad}+C^\ast_{\rm rad},
\qquad
\lambda_-^{\rm rad}=A_{\rm rad}-C^\ast_{\rm rad},
\]
entonces
\[
\boxed{
\frac{m(P)}{m_e}
= q_+^{-\frac{n_A+s_C}{2}}
q_-^{-\frac{n_A-s_C}{2}}.}
\]
En consecuencia, cada masa es un car\'acter multiplicativo de los dos canales espectrales del operador angular.
\end{teorema}

\begin{proof}
La identidad lineal
\[
n_AA+s_CC^\ast=\frac{n_A+s_C}{2}(A+C^\ast)+\frac{n_A-s_C}{2}(A-C^\ast)
\]
es inmediata. Pasando a radianes y usando
\[
q_+=e^{-(A_{\rm rad}+C^\ast_{\rm rad})},
\qquad
q_-=e^{-(A_{\rm rad}-C^\ast_{\rm rad})},
\]
se obtiene la f\'ormula de caracteres.
\end{proof}

\begin{observacion}[lectura MD]
La tabla de especies se convierte en una tabla de pesos sobre los dos idempotentes \(P_+\) y \(P_-\). La masa no es una asignaci\'on aislada: es un car\'acter del \`algebra split generada por el centro APP levantado a \((A,C^\ast)\).
\end{observacion}

\section{Teorema IX: extracci\'on nuclear de Barbero, Euler y Moonshine}

\begin{proposicion}[paso \(6\to8\) en el n\'ucleo]
La familia tr\'itica \(\mathcal B_\tau\) tiene canal diferencial
\[
\lambda_-(\tau)=6\tau-1.
\]
El salto fundamental entre hojas consecutivas es \(6\). La lectura del centro por filas da
\[
72=8\cdot9,
\qquad
27=3\cdot9.
\]
Por tanto, el n\'ucleo contiene simult\'aneamente un salto hex\'adico \(6\) y un sector oct\'adico \(8\). La lectura Barbero-HMT del paso
\[
90=6\binom62,
\qquad
120=8\binom62
\]
es la versi\'on areal de esta pareja nuclear.
\end{proposicion}

\begin{proposicion}[vac\'io diferencial y sombra \(-1/12\)]
En la familia \(\mathcal B_\tau\), el canal diferencial neutro es
\[
\lambda_-(0)=-1.
\]
Al distribuir este canal sobre la rueda dodecaf\'asica, se obtiene
\[
\boxed{-\frac1{12}.}
\]
\end{proposicion}

\begin{proposicion}[defecto comprimido y despliegue Moonshine]
La diferencia comprimida \(Q_\times-Q_+\) posee traza \(2\), suma \(6\) y despliegue global \(54\):
\[
\boxed{2\to6\to54.}
\]
El mismo \(54\) aparece como defecto global APP y como primer coeficiente positivo del Hauptmodul tri\'adico. Su amplificaci\'on HMT-Moonshine se expresa mediante identidades del tipo
\[
196884=54(5\cdot729+1).
\]
\end{proposicion}

\begin{observacion}[estado l\'ogico]
Las identidades aritm\'eticas de este apartado son exactas. La interpretaci\'on Barbero, Euler y Moonshine requiere los m\'odulos correspondientes: Barbero como lectura \(6\to8\) en \(90/120\), Euler-HMT como distribuci\'on dodecaf\'asica del vac\'io diferencial, y Moonshine como completaci\'on modular del defecto \(54\).
\end{observacion}

\section{Teorema X: tres firmas cuadr\'aticas del trit}

\begin{teorema}[tres algebras geom\'etricas]
La arquitectura HMT contiene tres endomorfismos de firma cuadr\'atica:
\[
A_W^2=-I,
\qquad
N^2=0,
\qquad
R^2=I.
\]
La familia formal
\[
\mathbb A_\tau=\mathbb K[J_\tau]/(J_\tau^2+\tau I),
\qquad \tau\in\{-1,0,+1\},
\]
organiza esas tres firmas:
\[
\tau=+1\Rightarrow J_\tau^2=-I,
\]
\[
\tau=0\Rightarrow J_\tau^2=0,
\]
\[
\tau=-1\Rightarrow J_\tau^2=I.
\]
\end{teorema}

\begin{proof}
La afirmaci\'on es formal: al imponer la relaci\'on \(J_\tau^2=-\tau I\) se obtienen las tres posibilidades. En el corpus, la hoja el\'iptica se realiza por la dualidad Witt \(A_W^2=-I\); la hoja parab\'olica por el operador Hensel \(N(x)=3x\) sobre \(\Z/9\Z\), para el cual \(N^2=0\); y la hoja hiperb\'olica por la involuci\'on two-way \(R^2=I\).
\end{proof}

\begin{observacion}[trit geom\'etrico]
El trit selecciona geometr\'ia: el\'iptica, parab\'olica o hiperb\'olica. En Mec\'anica Dimensional estas hojas se proyectan como curvatura positiva, contacto torsional y apertura hiperb\'olica.
\end{observacion}

\section{Arquitectura revisada del corpus}

La arquitectura que emerge de los teoremas anteriores se organiza en tres niveles.

\subsection*{N\'ucleo generativo}
\[
\mathfrak H_{\rm core}=\bigl(
\widehat{\mathbb T},
\mathcal A_9,
\mathsf K_{\rm TPK},
\mathcal K_{27},
\mathcal R_{12}^{\rm HMT},
\mathcal U_6,
\beta_3,
\mathcal G_9,
\mathcal L_{1000}
\bigr).
\]
Aqu\'i \(\mathcal G_9\) designa la ley de nueve puertas o monodrom\'ia de generaci\'on de regiones, y \(\mathcal L_{1000}\) la lectura cil\'indrica archimedeana.

\subsection*{Interfaz angular}
\[
\mathfrak I_{\rm ang}=\bigl(
\Bc,
A,
C^\ast,
\Bang,
q_+,
q_-
\bigr).
\]
La interfaz angular es la transducci\'on que conserva los proyectores del centro APP y sustituye el espectro entero \(9,5\) por el espectro archimedeano \(A+C^\ast,A-C^\ast\).

\subsection*{Proyecci\'on f\'isica}
\[
\mathfrak M_{\rm dim}=\bigl(
\mathfrak I_{\rm ang},
\Lambda,
\mathcal M,
\mathcal U_{\rm mix},
\mathcal C_{\rm phys}
\bigr).
\]
Aqu\'i \(\Lambda\) es el ledger de especies, \(\mathcal M\) el car\'acter de masas, \(\mathcal U_{\rm mix}\) el sistema de mezclas y \(\mathcal C_{\rm phys}\) el diccionario de carga, spin, color, sabor y conjugaci\'on.

La cadena formal queda:
\[
\boxed{\mathfrak H_{\rm core}\longrightarrow \mathfrak I_{\rm ang}\longrightarrow \mathfrak M_{\rm dim}.}
\]

\section{S\'intesis final}

El centro multiplicativo de la APP contiene el operador two-way fundamental de HMT. Su espectro entero \(9,5\) condensa el cierre multiplicativo y el centro aditivo. Su determinante \(45\) conserva la suma de la rueda. Su normalizaci\'on produce una forma Lorentz \(1+1\). Su familia tr\'itica extrae los tres reg\'imenes de geometr\'ia. Su compresi\'on Hensel despliega el defecto \(2\to6\to54\). Su lectura local \((72,27)\) comprime la red global \(3,6,9\). Su elevaci\'on archimedeana produce el operador angular \(AI+C^\ast R\). Sus canales propios generan \(q_+,q_-\). Las masas son caracteres de esos canales.

Por tanto:
\[
\boxed{\text{APP es el primer objeto donde coinciden aritm\'etica hologr\'afica, \`algebra hologr\'afica y geometr\'ia hologr\'afica.}}
\]

\[
\boxed{\text{El centro de APP es la bisagra espectral entre aritm\'etica, Lorentz y materia.}}
\]

\end{document}
